% GDZ PPN236018647,native537–538above50;HKUST538–539.
\clearpage\SourceLeaf{537}\ArithmeticHead{508}{L}{TAGEBUCH MIT ERLÄUTERUNGEN.}
\NumberedItem{44}
\indent Formula interpolationis elegans.
\par\noindent\hfill[1796] Nov. 25. G[ottingae]

\medskip\begingroup\small
\indent \textsc{Gauss} hat vielleicht hier die sogenannte \textsc{Lagrange}sche Interpolationsformel gemeint. Sie findet sich in \textsc{Lagranges} \textit{Leçons élémentaires sur les mathématiques} (Oeuvres de \textsc{Lagrange} VII, S. 286), die nach den Angaben in den Oeuvres zuerst in den zwei Ausgaben der Séances des Écoles normales, an III (1794—1795) erschienen sind und dann wieder im Journal de l’École Polytechnique, tome 2 (1812) abgedruckt wurden. Vorher war übrigens schon \textsc{Waring} zu dieser Formel gelangt. (Vgl. A. \textsc{von Braunmühl}, Bibliotheca mathematica (3) 2, 1904, S. 95 und 96). Bei \textsc{Gauss} steht die fragliche Formel in der aus seinem Nachlaß herausgegebenen Abhandlung \textit{Theoria interpolationis methodo nova tractata} (Werke III, S. 273), deren erster Entwurf nach \textsc{Schering} (ebenda, S. 328) aus dem Jahre 1805 zu stammen scheint. Auch in dem Aufsatz \textit{Über Interpolation} von J. F. \textsc{Encke} (Berliner Astronomisches Jahrbuch, 1830, Ges. math. u. astron. Abhandlungen, Berlin 1888, S. 4), der aus den bei \textsc{Gauss} im Jahre 1812 gehörten Vorlesungen hervorgegangen ist, spielt die Interpolationsformel eine grundlegende Rolle; endlich findet sie sich im art. 7. der Abhandlung \textit{Methodus nova integralium valores per approximationem inveniendi} (1814, Werke III, S. 172).
\par\noindent\hfill\textsc{Loewy.}
\par\endgroup

\NumberedItem{45}
\indent Incepi Expressionem
\[
1-\frac1{2^\omega}+\frac1{3^\omega}\ \cdotp\,\cdotp\,\cdotp
\]
\noindent in seriem transmutare secundum potestates ipsius \(\omega\) progredientem.
\par\noindent\hfill[1796] Nov. 26. G[ottingae]

\medskip\begingroup\small
\indent Mit der hier von \textsc{Gauss} betrachteten Reihe hat sich \textsc{Euler} in der Arbeit, Histoire de l’Académie Berlin, 17 (1761) 1768, Mémoires, S. 83 beschäftigt (vergl. E. \textsc{Landau}, Bibliotheca mathem. (3) 7, S. 69). Daß diese Reihe nur dann unbedingt konvergiert, wenn der reelle Teil von \(\omega\), \(\Re(\omega)\), größer ist als Eins, daß sie ferner, wie unsere Tagebuchnotiz angibt, in eine gewöhnliche Potenzreihe von \(\omega\) verwandelt werden kann und daß diese Potenzreihe dann beständig konvergent ist, folgt aus \textsc{Riemanns} Abhandlung, Monatsberichte der Berliner Akademie 1859, S. 671, Werke, 2. Aufl. 1892, S. 145. Nach neueren Untersuchungen ist die ursprüngliche Reihe für \(0<\Re(\omega)\leqq1\) bedingt konvergent.
\par\noindent\hfill\textsc{Schlesinger.}
\par\endgroup

\NumberedItem{46}
\indent Formulae trigonometricae per series expressae.
\par\noindent\hfill[1796] per Dec.

\NumberedItem{47}
\indent Differentiationes generalissimae.
\par\noindent\hfill[1796] Dec. 23.

\medskip\begingroup\small
\indent Es handelt sich jedenfalls um die Differentiation mit einem Index, der keine positive ganze Zahl ist.
\par\endgroup
\clearpage\SourceLeaf{538}\ArithmeticHead{509}{R}{1796 NOV. 25—1797 IAN. 7.}
\begingroup\small
\noindent Versuche, die sich auf solche Verallgemeinerungen beziehen, finden sich schon bei \textsc{Leibniz} in verschiedenen Briefen\footnote{Eine Zusammenstellung der einschlägigen Briefstellen gibt C. W. \textsc{Borchardt} in \textsc{Boncompagnis} Bulletino di bibliografia 2, 1869, S. 277, Ges. Werke, 1888, S. 486.}. \textsc{Gauss} hat wohl an \textsc{Euler} angeknüpft, siehe dessen Abhandlung \textit{De progressionibus transcendentibus etc.}, Comm. Acad. Petrop. 5 (1730/31) 1738, S. 36, besonders S. 55. Im Nachlaß von \textsc{Gauss} findet sich keine Spur dieser Untersuchungen.
\par\noindent\hfill\textsc{Schlesinger.}
\par\endgroup

\NumberedItem{48}
\indent Curvam parabolicam quadrare suscepi, cuius puncta quotcunque dantur.
\par\noindent\hfill[1796] Dec. 26.

\medskip\begingroup\small
\indent Hierher gehörige Formeln finden sich bei S. 13 des \textsc{Leiste}.
\par\noindent\hfill\textsc{Klein.}
\par\endgroup

\NumberedItem{49}
\indent Demonstrationem genuinam theorematis \textsc{La Grang}iani detexi.
\par\noindent\hfill[1796] Dec. 27.

\medskip\begingroup\small
\indent Aus dem oben S. 429 abgedruckten Briefe von \textsc{Gauss} an C. F. \textsc{Hindenburg} vom 8. Oktober 1799 geht hervor, daß \textsc{Gauss} etwa im April 1797 einen Aufsatz mit seinem Beweise des \textsc{Lagrange}schen Lehrsatzes durch \textsc{Kaestner} an \textsc{Hindenburg} hat senden lassen, damit dieser den Aufsatz in seinem Archiv der reinen und angew. Mathematik veröffentliche. Die Veröffentlichung ist unterblieben und, wie \textsc{Sartorius v. Waltershausen}\textsuperscript{*)} berichtet, »kam das eingesandte Manuskript nie wieder zum Vorschein«. Es befindet sich jedoch im Gaussarchiv eine Handschrift dieses Beweises, die Werke VIII, S. 76 abgedruckt ist. Eine etwas andere Darstellung desselben Beweises mit weniger Text ist im \textsc{Leiste} neben den Druckseiten 10—12 aufgezeichnet. Vergl. auch die Nr. 86 vom Mai 1798.
\begingroup\renewcommand{\thefootnote}{**)}\footnotetext{\textit{Gauss zum Gedächtniss}, 1856, S. 22; einige irrtümliche Angaben, die sich an dieser Stelle finden, sind schon oben S. 444 berichtigt worden.}\endgroup
\par\noindent\hfill\textsc{Klein. Schlesinger.}
\par\endgroup
